\documentclass[a4paper,12pt]{article}

\documentclass[a4paper,12pt]{article}

\usepackage{indentfirst}
\usepackage{cite}
\usepackage[utf8]{inputenc}
\usepackage[T1]{polski}
\usepackage{helvet}
\usepackage{graphicx}
\usepackage{svg}
\usepackage{color}
\usepackage{geometry}
\usepackage{float}
\usepackage{multirow}
\usepackage[hidelinks]{hyperref}
\usepackage{caption}
% dodaj bibliografię do spisu treści
\usepackage[nottoc]{tocbibind}
% unikaj pojedynczych linii na początku/ końcu stron
\usepackage[defaultlines=4,all]{nowidow}
\usepackage{subcaption}
\usepackage{amsmath}
\usepackage{amsfonts}

% set default figure placement to !htbp
\makeatletter
\def\fps@figure{!htbp}
\makeatother


\newcommand{\setSubtitle}[1]{
    \newcommand{\subtitle}{#1}
    }

\newcommand{\labdate}{\today}
\newcommand{\temat}{Echo}

\begin{document}
    
% Please add the following required packages to your document preamble:
% \usepackage{multirow}
% \usepackage{graphicx}

\begin{table}[]
    \resizebox{\textwidth}{!}{%
    \begin{tabular}{|c|c|c|c|}
    \hline
    \multicolumn{2}{|c|}{\multirow{2}{*}{\hspace{0.5cm} \includegraphics[height=2.3cm]{img/logoAGH} \hspace{0.5cm} }}                                                                    & \multicolumn{2}{c|}{\textbf{\begin{tabular}[c]{@{}c@{}} Akademia Górniczo-Hutnicza w Krakowie\\ Wydział Inżynierii Materiałowe i Ceramiki  \vspace{0.5cm} \end{tabular}}}                                     \\ \cline{3-4} 
    \multicolumn{2}{|c|}{}                                                                                       & \multicolumn{2}{c|}{\begin{tabular}[c]{@{}c@{}}Aleksandra Barcz\\ Alicja Kotwica  \vspace{0.5cm} \end{tabular}}                                                                                              \\ \hline
    \multicolumn{4}{|c|}{\textbf{Metody Badań Materiałów}}                                                                                                                                                                                                                                                     \\ \hline
    \multicolumn{4}{|c|}{\textbf{\temat}}                                                                                                                                                                                                                                                                            \\ \hline
    \begin{tabular}[c]{@{}c@{}}Rok Akademicki:\\ 2020/2021\end{tabular}         & \multicolumn{2}{c|}{\begin{tabular}[c]{@{}c@{}}Rok studiów:\\ IV\end{tabular}}                                  & \begin{tabular}[c]{@{}c@{}}Kierunek:\\ Ceramika\end{tabular}                                                       \\ \hline
    \begin{tabular}[c]{@{}c@{}}Data wykonania ćwiczenia\\ \labdate \end{tabular} & \multicolumn{2}{c|}{\begin{tabular}[c]{@{}c@{}}Data oddania sprawozdania:\\ \today \end{tabular}} & \begin{tabular}[c]{@{}c@{}}Prowadzący:\\ Dr. inż. Beata Macherzyńska\end{tabular}                                  \\ \hline
    \end{tabular}%
    }
    \end{table}


\section{Cel ćwiczenia}

\section{Wprowadzenie}

\section{Aparatura i materiały}

\section{Wyniki i obliczenia}

\subsection{Udział objętościowy}

$$V_V=\cfrac{\sum _i P_{ik}}{P_k}$$

Gdzie:

$V_V$ - udział objętościowy,

$P_{ik}$ - suma punktów w fazie,

$P_k$ - suma wszystkich punktów.

\subsection{Rozwinięcie powierzchni $S_V$}

$$S_V=\cfrac{\sum ^n_{i=1} S_i}{V}$$

Gdzie:

$S_V$ - rozwinięcie powierzchni,

$S_i$ - suma powierzchni,

$V$ - objętość, jaką zajmują ziarna.

$$\overline{P}_L=\cfrac{p}{k\cdot l}\sum^k_{i=1}n_i$$

Gdzie:

$\overline{P}_L$ - średnia liczba przecięć na jednostkę długości 

$p$ - powiększenie

$k$ - liczba pomiarów

$l$ - długość

$n_i$ - liczba punktów przecięcia z granicą faz

\subsection{Wielkość ziaren}

$$N_A=\cfrac{p^2}{a\cdot b}(z+0.5w+0.25u)$$

Gdzie:

$N_A$ - liczba ziaren mierzonych na jednostkowej powierzchni zgładu,

$p$ - powiększenie mikroskopu,

$a$ i $b$ - długość boków prostokąta,

$z$ - liczba ziaren leżących całkowicie wewnątrz prostokąta,

$w$ - liczba ziaren przeciętych przez brzegi prostokąta,

$u$ - liczba ziaren leżących w narożach prostokąta.


\section{Podsumowanie i wnioski}

\section{Załączniki wyników}

\end{document}
